% =====================================================================
%  8. Cronograma de execução
% =====================================================================
\section{Cronograma de execução}
\label{sec:cronograma}

O programa tem sete anos de execução, de 2027 a 2033, precedidos de um ano
de preparação institucional em 2026 (\autoref{fig:cronograma}). Os prazos
seguem a planilha de custos do PSH-PR \cite{custo2026}, que é a referência
para as metas anuais (\autoref{tab:metas-anuais}). Os ``Anos 0 a 6'' usados
nos quadros de etapas do MOP (Apêndice~\ref{ap:quadros}) correspondem aos
anos de 2027 a 2033. A estruturação institucional e social começa primeiro e
se estende por todo o período, porque o modelo de gestão precisa acompanhar os
sistemas até o encerramento. O trabalho técnico-social ocorre entre 2027 e
2030, e a implantação física, entre 2027 e 2033. Os sistemas do Programa
Sanepar Rural são entregues até 2031. Os do IAT e o esgotamento das
comunidades seguem de 2028 a 2033, e o esgotamento da área do IDR-Paraná, de
2029 a 2033.

\begin{figure}[htbp]
  \centering
  \caption{Cronograma de execução do Subcomponente 3.2.b}
  \label{fig:cronograma}
  % Cronograma (diagrama de Gantt) do Subcomponente 3.2.b
% Colunas: Ano 0 (coluna 1) a Ano 6 (coluna 7), conforme o MOP.
\begin{ganttchart}[
    hgrid,
    vgrid={*1{draw=gray!40, dotted}},
    x unit=1.15cm,
    y unit title=0.55cm,
    y unit chart=0.62cm,
    title/.append style={fill=azulclaro, draw=azulpsh},
    title label font=\scriptsize\bfseries,
    bar/.append style={fill=azulpsh!75, draw=azulpsh},
    bar height=0.55,
    bar label font=\scriptsize,
    bar label node/.append style={align=right, text width=6.2cm},
    group/.append style={fill=azulpsh, draw=azulpsh},
    group label font=\scriptsize\bfseries,
    group label node/.append style={align=right, text width=6.2cm},
    milestone/.append style={fill=orange, draw=orange!70!black},
    milestone label font=\scriptsize,
    milestone label node/.append style={align=right, text width=6.2cm}
  ]{1}{7}
  \gantttitle{Ano de execução}{7} \\
  \gantttitle{0}{1} \gantttitle{1}{1} \gantttitle{2}{1} \gantttitle{3}{1}
  \gantttitle{4}{1} \gantttitle{5}{1} \gantttitle{6}{1} \\

  \ganttgroup{I. Estruturação institucional e social}{1}{7} \\
  \ganttbar{3.2.2.1 Modelos de gestão}{1}{7} \\
  \ganttbar{3.2.2.2 Trabalho técnico-social}{2}{5} \\

  \ganttgroup{II. Implantação física}{2}{6} \\
  \ganttbar{3.2.2.3 Abastecimento de água}{3}{6} \\
  \ganttbar{3.2.2.4 Esgotamento -- comunidades}{3}{6} \\
  \ganttbar{3.2.2.5 Esgotamento -- área IDR-Paraná}{3}{6} \\
  \ganttbar{3.2.2.6 Sanepar Rural}{2}{6} \\

  \ganttgroup{III. Monitoramento e encerramento}{3}{7} \\
  \ganttbar{Supervisão técnica e monitoramento}{3}{7} \\
  \ganttbar{Relatórios finais e encerramento}{6}{7} \\
  \ganttmilestone{Termo de Recebimento Definitivo}{7}
\end{ganttchart}

  \fonte{Adaptado de \citeonline{custo2026} e do POA 2026--2027.}
\end{figure}

O Plano Operativo Anual (POA) detalha, mês a mês, o que acontecerá em 2026 e
2027 (Anexo~\ref{an:cronograma}). O ano de 2026 será dedicado à preparação.
A equipe interna do IAT fará o planejamento integrado, a seleção e a adesão
dos municípios, a assinatura dos Termos de Cooperação, a definição do modelo
de gestão piloto e o estudo de alternativas para o esgotamento. No segundo
semestre de 2026, começarão a ser elaborados os primeiros termos de
referência. Em 2027, as contratações estarão em licitação, e as primeiras
execuções estão previstas para o fim do ano. A mobilização social e os
estudos do poço começarão em outubro, a perfuração em novembro e os projetos
das redes em dezembro de 2027. As demais contratações não iniciam a execução em
2027. As aquisições de materiais para as redes de água e para o esgotamento,
inclusive na frente do IDR-Paraná, e a supervisão das obras estarão em
licitação até dezembro de 2027. O segundo lote de materiais de água e a
infraestrutura do modelo de gestão só terão os TR iniciados no fim de 2027.

\pendencia{Atualizar no MOP todas as referências a ``seis anos'' e aos
``Anos 0 a 6'' para o prazo de sete anos (2027 a 2033).}
